\sloppy

% Overleaf compatibility switch.
% true  - skip external title PDF and use listings instead of minted
% false - use local mode with minted and Title.pdf
\newif\ifoverleafmode
\overleafmodetrue

% Настройки стиля ГОСТ 7-32
% Для начала определяем, хотим мы или нет, чтобы рисунки и таблицы нумеровались в пределах раздела, или нам нужна сквозная нумерация.
% \EqInChapter % формулы будут нумероваться в пределах раздела
% \TableInChapter % таблицы будут нумероваться в пределах раздела
% \PicInChapter % рисунки будут нумероваться в пределах раздела

% Добавляем гипертекстовое оглавление в PDF
\usepackage[
bookmarks=true, colorlinks=true, unicode=true,
urlcolor=black,linkcolor=black, anchorcolor=black,
citecolor=black, menucolor=black, filecolor=black,
]{hyperref}

\urlstyle{same}

\usepackage{pdfpages}
\usepackage[all]{hypcap}
\usepackage{needspace}

% Изменение начертания шрифта --- после чего выглядит таймсоподобно.
% apt-get install scalable-cyrfonts-tex

% \IfFileExists{cyrtimes.sty}
%     {
        % \usepackage{cyrtimespatched}
%     }
%     {
%         % А если Times нету, то будет CM...
%     }

\usepackage{graphicx}   % Пакет для включения рисунков
\graphicspath{{images/}}

% С такими оно полями оно работает по-умолчанию:
% \RequirePackage[left=20mm,right=10mm,top=20mm,bottom=20mm,headsep=0pt]{geometry}
% Если вас тошнит от поля в 10мм --- увеличивайте до 20-ти, ну и про переплёт не забывайте:

% Пакет Tikz
\usepackage{tikz}
\usetikzlibrary{arrows,positioning,shadows}

% Произвольная нумерация списков.
\usepackage{enumerate}

% ячейки в несколько строчек
\usepackage{multirow}

% itemize внутри tabular
\usepackage{paralist,array}

% Таблицы
\usepackage{array} % расширенные возможности для работы с таблицами
\usepackage{tabularx} % автоматический подбор ширины столбцов
\usepackage{dcolumn} % выравнивание чисел по разделителю

\usepackage{setspace} % полуторный интервал
\setstretch{1.43}

\newenvironment{compactlist}{%
    \renewcommand{\labelenumi}{\hspace{13mm}\arabic{enumi}.}
    \setlength{\topsep}{0pt}%
    \setlength{\partopsep}{0pt}%
    \setlength{\itemsep}{0pt}%
    \setlength{\parsep}{0pt}%
    \begin{enumerate}%
}{%
    \end{enumerate}%
}

\makeatletter
\renewcommand*\l@section{\@dottedtocline{1}{0em}{2.3em}}
\renewcommand*\l@subsection{\@dottedtocline{2}{0em}{3.2em}}
\renewcommand*\l@subsubsection{\@dottedtocline{3}{0em}{4.1em}}
\renewcommand\@dotsep{2.2}
\makeatother


\usepackage{caption}
\usepackage{ragged2e}

\DeclareCaptionJustification{centerednohyphen}{\Centering\hyphenpenalty=10000\exhyphenpenalty=10000}
\DeclareCaptionLabelSeparator{shortdash}{ \textendash\ }

\usepackage{caption}

\usepackage{float}
\setlength{\floatsep}{0pt}
\setlength{\textfloatsep}{0pt}
\setlength{\intextsep}{7pt}

\captionsetup{aboveskip=7pt,belowskip=0pt,justification=centerednohyphen,singlelinecheck=false,labelsep=shortdash}
\captionsetup[table]{justification=RaggedRight,singlelinecheck=false}
\captionsetup[figure]{justification=centerednohyphen,singlelinecheck=false,font=normalsize,labelfont=normalsize}
\captionsetup[listing]{justification=RaggedRight,singlelinecheck=false, position=top}
\captionsetup[listing]{position=top}

\usepackage{chngcntr}

\usepackage{mdframed}
\usepackage{xcolor}
\usepackage{etoolbox}

\mdfsetup{
linecolor=black,
linewidth=1pt,
roundcorner=5pt,
innerleftmargin=5pt,
innerrightmargin=5pt,
innertopmargin=5pt,
innerbottommargin=5pt,
skipabove=0pt, % Reduces space above the frame
skipbelow=0pt, % Reduces space below the frame
leftmargin=0pt, % Reduces left outer margin
rightmargin=0pt % Reduces right outer margin
}

\ifoverleafmode
  \usepackage{listings}
  \usepackage{newfloat}
  \DeclareFloatingEnvironment[
    fileext=lol,
    listname={Листинги},
    name=Листинг
  ]{listing}
  \lstset{
      basicstyle=\ttfamily\small,
      numbers=left,
      numberstyle=\scriptsize\color{gray},
      numbersep=8pt,
      breaklines=true,
      breakatwhitespace=false,
      columns=fullflexible,
      keepspaces=true,
      showstringspaces=false,
      xleftmargin=2.4em,
      framexleftmargin=2.0em,
      tabsize=2,
      backgroundcolor=\color{black!3},
      keywordstyle=\color{blue!70!black}\bfseries,
      commentstyle=\color{green!45!black}\itshape,
      stringstyle=\color{red!65!black}
  }
  \lstnewenvironment{minted}[2][]{
    \lstset{
      language=#2,
      #1
    }
  }{}
  \newcommand{\inputminted}[3][]{\lstinputlisting[language=#2,#1]{#3}}
\else
  \usepackage[chapter]{minted}
  \usemintedstyle{friendly}

  % Configure captions for minted's listing environment
  \captionsetup[listing]{
      position=top,
      justification=RaggedRight,
      singlelinecheck=false
  }

  % Alternative: Force all float captions to top
  \floatstyle{plaintop}
  \restylefloat{listing}

  \setminted{
      linenos=true,           % включить номера строк
      numberblanklines=false, % не нумеровать пустые строки
      numbersep=5pt,          % расстояние между номерами строк и кодом
      xleftmargin=1em,        % отступ слева, чтобы номера не вылезали за границы
      fontsize=\small,
      breakanywhere=true,
      breakautoindent=true,
      breaklines=true,
      bgcolor=black!3,
  }

  \renewcommand{\listingscaption}{Листинг}
  \renewcommand{\listoflistingscaption}{Листинги}
\fi


\makeatletter
\let\oldtableofcontents\tableofcontents
\renewcommand{\tableofcontents}{%
  \begingroup
  \hyphenpenalty=10000
  \exhyphenpenalty=10000
  \oldtableofcontents
  \endgroup
}
\makeatother

\makeatletter
\def\@hangfrom#1{\setbox\@tempboxa\hbox{{#1}}%
      \hangindent 0pt%\wd\@tempboxa
      \noindent\box\@tempboxa}
\makeatother

\ifoverleafmode\else
\AtEndEnvironment{listing}{\vspace{-4pt}}
\fi
\AtEndEnvironment{table}{\vspace{0pt}}
\AtEndEnvironment{center}{\vspace{0pt}}
\AtEndEnvironment{figure}{\vspace{0pt}}


\usepackage{fontspec}

% Указываем шрифты для основного текста и для кириллицы:
\setmainfont{Times New Roman} % если установлен в системе
\newfontfamily\cyrillicfont{Times New Roman} % для кириллического текста

% Убирает точку после номера главы
\makeatletter
\renewcommand{\thesection}{\thechapter.\arabic{section}}
\renewcommand{\thesubsection}{\thesection.\arabic{subsection}}
\renewcommand{\@seccntformat}[1]{%
  {\fontsize{14pt}{21pt}\selectfont\bfseries \csname the#1\endcsname}\hspace{1em}%
}
\makeatother

% Keep normal spacing after section titles.
\setlength\GostAfterTitleSkip{3.2ex}

% Remove bottom spacing for second-level headings (\section).
\makeatletter
\renewcommand\section{%
 \@afterindenttrue
 \@startsection{section}{1}%
 {\GostTitleIndent}{\GostBeforTitleSkip}{0.6ex}%
 {\GostTitleStyle}%
}
\makeatother
